% Part: incompleteness
% Chapter: introduction
% Section: historical-background

\documentclass[../../../include/open-logic-section]{subfiles}

\begin{document}

\olfileid[id]{inc}{int}{bgr}

\olsection{Latar Belakang Historis}

Dalam bagian ini, kita akan membahas secara singkat perkembangan historis
yang membantu menempatkan teorema-teorema ketaklengkapan dalam konteksnya.
Secara khusus, kita akan memberikan ikhtisar yang sangat garis besar tentang
sejarah logika matematika, lalu menyampaikan beberapa hal mengenai sejarah
landasan matematika.

\begin{digress}
  Frasa ``logika matematika'' bersifat ambigu. Kata ``matematika'' dapat
  dipahami sebagai penjelas pokok bahasan, seperti dalam ``logika
  matematika,'' yang menunjuk pada prinsip-prinsip penalaran matematis; atau
  sebagai penjelas metode, seperti dalam ``matematika tentang logika,'' yang
  menunjuk pada kajian matematis atas prinsip-prinsip penalaran. Uraian
  berikut melibatkan logika matematika dalam kedua arti tersebut, sering kali
  secara bersamaan.
\end{digress}

Kajian logika pada dasarnya bermula dari Aristoteles, yang hidup sekitar
384--322 \textsc{sm}. Karyanya \emph{Categories}, \emph{Prior Analytics},
dan \emph{Posterior Analytics} memuat kajian sistematis tentang
prinsip-prinsip penalaran ilmiah, termasuk kajian silogisme yang menyeluruh
dan sistematis.

Logika Aristoteles mendominasi filsafat skolastik sepanjang Abad Pertengahan;
bahkan hingga abad kedelapan belas, Kant berpendapat bahwa logika Aristoteles
sudah sempurna dan tidak memerlukan revisi. Namun, teori silogisme terlalu
terbatas untuk memodelkan lebih daripada aspek-aspek paling dangkal dari
penalaran matematis. Seabad sebelumnya, Leibniz, sezaman dengan Newton,
membayangkan suatu ``kalkulus'' lengkap bagi penalaran logis dan mengambil
beberapa langkah awal untuk merancang kalkulus tersebut, yang pada dasarnya
mendeskripsikan suatu versi logika proposisional.

Abad kesembilan belas merupakan titik balik bagi logika. Pada 1854 George
Boole menulis \emph{The Laws of Thought}, dengan kajian aljabar yang mendalam
atas logika proposisional dan tidak jauh dari penyajian modern. Pada 1879
Gottlob Frege menerbitkan \emph{Begriffsschrift} (Tulisan Konsep), yang
memperluas logika proposisional dengan kuantor dan relasi sehingga mencakup
logika orde pertama. Bahkan, sistem logika Frege juga mencakup logika orde
tinggi dan lebih banyak lagi. Dalam \emph{Basic Laws of Arithmetic}, Frege
berusaha menunjukkan bahwa seluruh aritmetika dapat diturunkan dalam
Begriffsschrift-nya dari asumsi-asumsi yang murni logis. Sayangnya,
asumsi-asumsi itu ternyata tak konsisten, sebagaimana ditunjukkan Russell pada
1902. Namun, jika aksioma yang tak konsisten itu disisihkan, Frege kurang
lebih menciptakan logika modern seorang diri---sebuah pencapaian yang
menakjubkan. Logika kuantifikasional juga dikembangkan secara independen oleh
para pemikir berorientasi aljabar setelah Boole, termasuk Peirce dan
Schr\"oder.

Sekarang mari beralih ke perkembangan dalam landasan matematika. Tentu saja,
karena logika memainkan peran penting dalam matematika, terdapat banyak
interaksi dengan perkembangan yang baru dijelaskan. Sebagai contoh, Frege
mengembangkan logikanya dengan tujuan eksplisit untuk menunjukkan bahwa
seluruh matematika dapat didasarkan hanya pada kerangka logisnya; secara
khusus, ia ingin menunjukkan bahwa matematika terdiri atas kebenaran
\emph{analitik} a priori, bukan kebenaran \emph{sintetik} a priori seperti
yang dipertahankan Kant.

Banyak orang menganggap kelahiran matematika dalam arti sebenarnya terjadi
bersama bangsa Yunani. \emph{Elements} karya Euklides, yang ditulis sekitar
300 SM, sudah merupakan contoh matang matematika Yunani dengan penekanannya
pada keketatan dan ketepatan. Definisi dan pembuktian dalam \emph{Elements}
bertahan kurang lebih utuh dalam buku teks geometri sekolah menengah hingga
kini (sejauh geometri masih diajarkan di sekolah menengah). Model penalaran
matematis ini dipandang sebagai paradigma argumentasi ketat, bukan hanya dalam
matematika tetapi juga dalam cabang-cabang filsafat. (Spinoza bahkan
menyajikan argumen moral dan keagamaan dalam gaya Euklides, yang terasa aneh
ketika dibaca!)

Kalkulus diciptakan oleh Newton dan Leibniz pada abad ketujuh belas. (Sengketa
keras tentang siapa yang lebih dahulu berlangsung selama berabad-abad, tetapi
kebanyakan sarjana masa kini berpendapat bahwa kedua perkembangan tersebut
sebagian besar berlangsung secara independen.) Kalkulus melibatkan penalaran,
misalnya, tentang jumlah tak hingga dari besaran infinitesimal;
ciri-ciri ini memicu kritik Uskup Berkeley, yang berpendapat bahwa kepercayaan
kepada Tuhan tidak kurang rasional daripada matematika pada masanya. Metode
kalkulus digunakan secara luas pada abad kedelapan belas, misalnya oleh
Leonhard Euler, yang menggunakan perhitungan yang melibatkan jumlah tak hingga
dengan hasil yang mencolok.

Pada abad kesembilan belas, matematikawan berusaha menjawab kritik Berkeley
dengan meletakkan kalkulus di atas landasan yang lebih kukuh. Upaya Cauchy,
Weierstrass, Bolzano, dan lainnya menghasilkan definisi modern tentang limit,
kekontinuan, diferensiasi, dan integrasi dalam istilah ``epsilon dan delta,''
dengan kata lain tanpa mengacu pada infinitesimal. Pada masa selanjutnya
dalam abad itu, matematikawan mencoba melangkah lebih jauh dan menjelaskan
semua aspek kalkulus, termasuk bilangan real itu sendiri, dalam istilah
bilangan asli. (Kronecker: ``Tuhan menciptakan bilangan bulat; segala yang
lain adalah karya manusia.'') Pada 1872, Dedekind menulis ``Continuity and
the Irrational Numbers,'' tempat ia menunjukkan cara ``membangun'' bilangan
real sebagai himpunan bilangan rasional (yang, sebagaimana diketahui, dapat
dipandang sebagai pasangan bilangan asli); pada 1888 ia menulis ``Was sind
und was sollen die Zahlen'' (kira-kira, ``Apakah bilangan asli itu dan
bagaimana seharusnya?''), yang bertujuan menjelaskan bilangan asli dalam
istilah yang murni ``logis.'' Pada 1887 Kronecker menulis ``\"Uber den
Zahlbegriff'' (``Tentang konsep bilangan''), tempat ia membahas representasi
semua objek matematis dalam istilah bilangan bulat; pada 1889 Giuseppe Peano
memberikan aksioma formal simbolis bagi bilangan asli.

Akhir abad kesembilan belas juga membawa keberanian baru dalam menangani yang
tak hingga. Sebelumnya, objek dan struktur infinit (seperti himpunan bilangan
asli) diperlakukan dengan sangat hati-hati; ``tak hingga banyaknya'' dipahami
sebagai ``sebanyak yang diinginkan,'' dan ``mendekati dalam limit'' dipahami
sebagai ``sedekat yang diinginkan.'' Namun, Georg Cantor menunjukkan bahwa
yang tak hingga dapat diperlakukan apa adanya. Karya Cantor, Dedekind, dan
lainnya membantu memperkenalkan pemahaman teori-himpunan umum tentang
matematika yang sekarang diterima luas.

Hal ini membawa kita ke perkembangan logika dan landasan pada abad kedua
puluh. Pada 1902 Russell menemukan paradoks dalam sistem logis Frege. Pada
1904 Zermelo membuktikan prinsip pengurutan-baik Cantor dengan menggunakan
apa yang disebut ``aksioma pilihan''; keabsahan aksioma ini memicu banyak
perdebatan. Antara 1910 dan 1913 terbit tiga jilid \emph{Principia
Mathematica} karya Russell dan Whitehead, yang memperluas program Frege untuk
mendasarkan matematika pada landasan logis. Sayangnya, Russell dan Whitehead
terpaksa menerima dua prinsip yang tampak sulit dibenarkan sebagai prinsip
yang murni logis: aksioma ketakhinggaan dan aksioma ``ketereduksian.'' Pada
dekade 1900-an, Poincar\'e mengkritik penggunaan ``definisi impredikatif''
dalam matematika, dan pada dekade 1910-an Brouwer mulai mengusulkan agar
seluruh matematika dilandaskan kembali pada basis ``intuisionistik,'' yang
menghindari hukum tiada jalan tengah ($!A \lor \lnot !A$).

Sungguh masa yang aneh!{} Program mereduksi seluruh matematika menjadi logika
kini disebut ``logisisme'' dan umumnya dipandang gagal akibat kesulitan yang
disebutkan di atas. Program mengembangkan matematika dalam istilah konstruksi
mental intuisionistik disebut ``intuisionisme'' dan dipandang menerapkan
pembatasan yang terlalu berat pada matematika sehari-hari. Sekitar pergantian
abad, David Hilbert, salah satu matematikawan paling berpengaruh sepanjang
masa, merupakan pendukung kuat metode baru dan abstrak yang diperkenalkan
Cantor dan Dedekind: ``tidak seorang pun akan mengusir kita dari surga yang
diciptakan Cantor bagi kita.'' Pada saat yang sama, ia peka terhadap kritik
landasan atas metode baru ini (yang, secara ganjil, sekarang disebut
``klasik''). Ia mengusulkan cara untuk memperoleh keduanya:
\begin{enumerate}
\item Representasikan metode klasik dengan aksioma dan aturan formal;
  representasikan persoalan matematis sebagai !!{formula}s dalam suatu sistem
  aksiomatik.
\item Gunakan metode ``finitistik'' yang aman untuk membuktikan bahwa sistem
  deduktif formal ini konsisten.
\end{enumerate}

Karya Hilbert sudah membawa kemajuan besar menuju tujuan pertama. Pada 1899,
ia telah melakukannya untuk geometri dalam bukunya yang terkenal,
\emph{Foundations of Geometry}. Pada tahun-tahun berikutnya, ia dan sejumlah
murid serta koleganya mengerjakan bidang matematika lain untuk melakukan apa
yang telah dilakukan Hilbert bagi geometri. Hilbert sendiri memberikan
sistem aksioma bagi aritmetika dan analisis. Zermelo memberikan aksiomatisasi
teori himpunan, yang dikembangkan lebih lanjut oleh Fraenkel, Skolem, von
Neumann, dan lainnya. Menjelang pertengahan dekade 1920-an, terdapat dua
pendekatan yang sama-sama mengklaim sebagai aksiomatisasi ``seluruh''
matematika: \emph{Principia Mathematica} karya Russell dan Whitehead, serta
apa yang kemudian dikenal sebagai teori himpunan Zermelo--Fraenkel.

Pada 1921, Hilbert memulai proyek penelitian untuk mencapai tujuan
membuktikan konsistensi sistem-sistem tersebut. Ia dibantu oleh beberapa
muridnya, khususnya Bernays, Ackermann, dan kemudian Gentzen. Gagasan dasarnya
ialah mengubah pertanyaan tentang kemungkinan adanya !!{derivation} suatu
inkonsistensi dalam matematika menjadi persoalan kombinatorial mengenai
barisan simbol yang mungkin, yakni barisan kalimat yang memenuhi kriteria
sebagai !!{derivation} yang benar atas, misalnya, $!A \land \lnot !A$ dari
aksioma sistem aksioma bagi aritmetika, analisis, atau teori himpunan. Bukti
bahwa barisan simbol semacam itu mustahil ada---karena bukti itu sendiri
merupakan bukti matematis---akan dapat diformalkan dalam sistem aksiomatik
tersebut. Dengan kata lain, akan ada suatu kalimat $\OCon$ yang menyatakan
bahwa, misalnya, aritmetika konsisten. Lebih jauh, kalimat ini seharusnya dapat
dibuktikan dalam sistem yang dimaksud, khususnya jika pembuktiannya hanya
memerlukan sarana ``finitistik'' yang sangat terbatas.

Tujuan kedua, yaitu agar sistem aksioma yang dikembangkan menyelesaikan setiap
persoalan matematis, dapat dibuat tepat dalam dua cara. Cara pertama ialah
merumuskannya sebagai berikut: untuk setiap kalimat~$!A$ dalam bahasa suatu
sistem aksioma bagi matematika, $!A$ atau $\lnot !A$ dapat dibuktikan dari
aksioma-aksiomanya. Jika hal ini benar, tidak akan ada kalimat yang tidak dapat
dibuktikan maupun disangkal berdasarkan aksioma, dan tidak ada pertanyaan yang
tidak diselesaikan oleh aksioma. Sistem aksioma dengan sifat ini disebut
\emph{lengkap}. Tentu saja, untuk kalimat tertentu mungkin tetap sulit
menentukan alternatif mana yang berlaku. Namun, pada prinsipnya seharusnya ada
metode untuk menentukannya. Bahkan, untuk sistem aksioma dan !!{derivation}
yang dipertimbangkan Hilbert, kelengkapan akan mengimplikasikan adanya metode
semacam itu---meskipun Hilbert tidak menyadarinya. Cara kedua untuk memahami
pertanyaan tersebut ialah persyaratan yang lebih kuat: adanya metode mekanis
dan komputasional yang, untuk suatu kalimat~$!A$, menentukan apakah kalimat itu
dapat diturunkan dari aksioma atau tidak.

Pada 1931, G\"odel membuktikan dua ``teorema ketaklengkapan,'' yang
menunjukkan bahwa program ini tidak dapat berhasil. Tidak ada sistem aksioma
efektif dan cukup kuat bagi aritmetika yang sekaligus konsisten dan lengkap. Secara
khusus, teori semacam itu tidak dapat membuktikan kalimat yang menyatakan
konsistensinya sendiri.

Hasil ini memberi pukulan mematikan bagi program Hilbert yang semula. Namun,
seperti yang sering terjadi dalam matematika, hasil tersebut juga membuka
jalur penelitian baru yang menarik. Jika tidak ada satu sistem formal yang
mencakup seluruh matematika, masuk akal untuk mengembangkan sistem yang
lingkupnya lebih terbatas dan menyelidiki apa yang dapat dibuktikan di
dalamnya. Masuk akal pula untuk mengembangkan metode pembuktian yang tidak
terlalu dibatasi guna menetapkan konsistensi sistem tersebut, serta menemukan
cara mengukur seberapa sulit konsistensinya dibuktikan. Karena G\"odel
menunjukkan bahwa (hampir) setiap sistem formal memiliki pertanyaan yang
tidak dapat diselesaikannya, masuk akal untuk mencari pertanyaan ``menarik''
yang tidak dapat diselesaikan suatu sistem formal dan menentukan seberapa kuat
sistem formal yang diperlukan untuk menyelesaikannya. Hingga kini, para ahli
logika terus meneliti pertanyaan tersebut dalam disiplin matematika baru,
yakni teori pembuktian.

\end{document}
